\section{Methodology}\label{sec:method}

In this section we present the proposed gradual pruning strategy utilizing the Iteratively Rewighted $l_1$ Minimization algorithm presented in Section
\ref{subsec:irl1}. Additionally, we describe two design choices and their motivations in the pruning process.

\subsection{IRL1 Regularization}

As explained in previous sections, utilizing a sparsity-enhancing regularization term in the cost function forces the optimization process
to find a sparser model than the initial one while fitting the data. The IRL1 algorithm enables us to approximate the most effective regularization term,
the $l_0$ norm penalty term using gradient descent optimization. Regarding weight pruning methods, applying a global threshold for pruning
facilitates sparser but accurate models due its flexibility in enhancing the sparsity of layers with different positions in the network unevenly.  
By applying the threshold multiple times, the initial dense model can be gradually compressed to a desired sparsity level.

Algorithm \ref{tab:pruning_algorithm} blends these two techniques into a single iterative pruning framework. The algorithm consists of
two phases. The first phase is a pre-training phase where the initial dense model is trained using the standard cross-entropy loss function for
$e$ epochs. After this, the resulting weights $\hat{W}$ - and optimizer specific buffers -  are saved as a checkpoint

The second phase is the main iterative pruning phase, where the model is gradually compressed to the target sparsity level $P$. It consists of two
sub-phases. The first one is defined as the regularization phase, where the model is trained using the cost function of Eq. \ref{eq:reg_cost} with 
the weighted $l_1$ regularization term defined in Eq. \ref{eq:wl1_term}. The penalty weights $A^{(\ell)}$ are updated after $n$ consecutive
epochs, according to Eq. \ref{eq:weight_update}.

To eliminate the hyperparameter related to the length of the regularization phase, an early-stopping criterion is applied based on
the difference in sparsity $\Delta p_t$ between two consecutive epochs. 
If the increase in sparsity is smaller than a predefined tolerance $\epsilon$ for $m$ consecutive epochs, then 
the regularization phase is stopped and the second sub-phase starts.

After updating the binary mask $Z$ based on the global threshold $\tau$ and the regularized weights $W$, a re-fitting iteration is applied
to let the remaining weights adapt to the new structure and recover accuracy. After the convergence of the latter training, 
the achieved sparsity is compared to the target sparsity $P$. If the current sparsity $p$ is below the target, Phase 2 is repeated.

As discussed in Section \ref{subsec:rewind}, rewinding the remaining weights after pruning to states from a dense training checkpoint can
lead to better performance than fine-tuning the pruned model from its current state. Therefore, the proposed algorithm
restores the weights to checkpoint $\hat{W}$ after each pruning step, while also applying the newly discovered mask $Z$. Note that this rewinding step
is one of the design choices whose benefits and drawback we evaluated during the experiments.

\begin{table}[H]
    \centering
    \renewcommand{\arraystretch}{1}
    \begin{tabular}{@{}ll@{}}
        \toprule
        \multicolumn{2}{l}{\textbf{Algorithm 1} Iterative Pruning Framework} \\
        \midrule
        \textbf{Input:}  & Training dataset $\mathcal{D}_{train}$, Validation dataset $\mathcal{D}_{val}$ \\
                         & Target sparsity $P$, Patience epochs $x$, Sparsity tolerance $E$ \\
                         & Global magnitude threshold $\tau$ \\
                         & Update interval $n$, \\
        \textbf{Output:} & Pruned weights $W^*$, Final binary mask $Z^*$ \\
        \midrule
        \multicolumn{2}{l}{\textbf{Phase 1: Dense Pretraining}} \\
        1: & Initialize dense weights $W^{(0)}$ \\
        2: & Train $W^{(0)}$ minimizing cross-entropy loss $\mathcal{L}_{CE}$ on $\mathcal{D}_{train}$ \\
        3: & At epoch $e$, save current weights to checkpoint $\hat{W}$ \\
        \midrule
        \multicolumn{2}{l}{\textbf{Phase 2: Iterative Pruning}} \\
        4: & Initialize current sparsity $p \gets 0$, binary mask $Z \gets \mathbf{1}$, and penalty $A \gets \mathbf{1}$ \\
        5: & \textbf{while} $p < P$ \textbf{do} \\
        \multicolumn{2}{l}{\quad \textit{Phase 2.1: Regularization}} \\
        6: & \quad Mini-batch training of $W$ on $\mathcal{D}_{train}$ minimizing $\mathcal{L} = \mathcal{L}_{CE} + \sum (A \odot \lvert W \rvert)$ \\
        7: & \quad Update penalty $A$ according to Equation 3 \\
        8: & \quad ............................. \\
        9: & \quad \textbf{if} $\Delta p_t = (p_t - p_{t-1}) < \epsilon$ for $x$ consecutive epochs \textbf{then} \\
        10:& \quad \quad \textbf{break} mini-batch training \\
        11:& \quad \textbf{end if} \\
        \multicolumn{2}{l}{\quad \textit{Phase 2.2: Mask Update \& Optimization}} \\
        12:& \quad $Z \gets (\lvert W \rvert > \tau)$ \\
        13:& \quad Restore and mask checkpoint weights: $W \gets \hat{W} \odot Z$ \\
        14:& \quad Optimize weights: $W \gets \arg\min_{W} \mathcal{L}_{CE}(W, \mathcal{D}_{train})$ \\
        15:& \quad $p \gets 1 - \frac{\|Z\|_0}{N}$ \quad \textit{(where $N$ is the total number of parameters)} \\
        16:& \textbf{end while} \\
        17:& \textbf{return} $W, Z$ \\
        \bottomrule
    \end{tabular}
    \caption{Algorithm: Iterative Pruning with Weighted $L_1$ Regularization}
    \label{tab:pruning_algorithm}
\end{table}


\subsection{Scheduling $\epsilon$ and $\lambda$}

Besides weight rewinding, the other design choices we focused on were the iterative decaying of $\epsilon$ and iterative increasing of $\lambda$ parameters.
The motivation behind this scheduling is to achieve a smoother transition from training on just the fidelity term to training 
on the cost function with the additional regularization term. By gradually increasing the pressure that regularization imposes on the weights,
we can avoid the optimization process being overwhelmed by the regularization term at the beginning of the regularization phase. 

To visualize the effect of introducing this smooth transition, we inspected the dynamics of the thresholded sparsity curve during the
gradual pruning process. The thresholded sparsity is defined as the percentage of weights that are below the global threshold $\tau$ at each epoch.

Figure \ref{fig:schedule} shows the different dynamics in the thresholded sparsity of the whole model among the aforementioned two regularization strategies.

\begin{figure}[H]
    \includegraphics[width=0.8\linewidth]{figs/curve.png}
    \centering
    \caption{Scheduling of $\epsilon$ and $\lambda$ during the training process.}
    \label{fig:schedule}
\end{figure}

We can conclude that continuously decaying $\epsilon$ and increasing $\lambda$ leads to a more stable optimization process. Based on 
the dynamics of the sparsity curve, scheduling these parameters achieves a smoother increase, while the discrete step 
makes the sparsity oscillate. The latter is due to gradients from the fidelity term being close to zero at the beginning of the regularization phase, 
as the weights are converged to a local minimum by the end of the re-fitting phase. Hence, in the first mini-batch update of the first epoch of
regularization the optimizer has a history of near-zero gradients so the optimizer
takes a large, unnormalized step in the direction of the regularization term gradient, while neglecting the fidelity term gradient and pushing
a large group of weights below the threshold. Then, in the next mini-batch update, the data loss gradients are no longer close to zero, 
and they pull the weights back up. By the end of the first regularization epoch, this oscillation is dampened by the optimizer's momentum, 
but it still leads to a significant jump in the thresholded sparsity. As the regularization phase continues, weights near the threshold are pushed 
below and above the threshold alternately, leading to the observed oscillation in the sparsity curve.

On the other hand, the scheduling strategy eliminates the effect of unnormalized gradient steps and achieves a slower but more
stable increase in sparsity.

If our assumption is that a smoother increase in sparsity indicates that regularization does not overwhelm the fidelity term of the loss function, and enables
the optimization process to find a better sparser architecture, then the scheduling strategy should achieve higher accuracy at the same sparsity level
than its discrete counterpart.

\subsection{Robustness Metrics}
To facilitate a broader comparison of the two strategies, we utilized two complementary metrics regarding model robustness. The first one
measures the smoothness of the input-output mapping the network implements.

Let $f(x; W) \in \mathbb{R}^C$ denote the neural network's output logits for a given input sample $x \in \mathbb{R}^D$, parameterized by the weights $W$. To evaluate the smoothness of the learned input-output mapping, we analyze the sensitivity of the network to local input perturbations.

We define the expected smoothness metric, $\bar{\mathcal{J}}_{\mathrm{smooth}}$, as the average Frobenius norm of the input Jacobian across the dataset $\mathcal{D}_{\mathrm{train}}$:
\begin{equation}
    \bar{\mathcal{J}}_{\mathrm{smooth}} = \mathbb{E}_{x \sim \mathcal{D}_{\mathrm{train}}} \left[ \left\| \frac{\partial f(x; W)}{\partial x} \right\|_F \right]
\label{eq:jacobian_norm}
\end{equation}

A lower value of $\bar{\mathcal{J}}_{\mathrm{smooth}}$ implies a lower upper bound on the empirical Lipschitz constant, indicating that the network implements
a smoother, more robust mapping from the input space to the output logits.

The second metric aims to quantify how the learned weight matrices of the network amplify noise or small changes in the input.

Let $W$ be the weight matrix of a given layer. The effective condition number, denoted as $\kappa_{\mathrm{eff}}(W)$, is defined as the ratio of the
largest singular value to the smallest non-zero singular value of $W$:
\begin{equation}
    \kappa_{\mathrm{eff}}(W) = \frac{\sigma_{\max}(W)}{\sigma_{\min}^{+}(W)}
\label{eq:conditionnumber}
\end{equation}
where $\sigma_{\max}(W)$ is the maximum singular value, and $\sigma_{\min}^{+}(W)$ represents the smallest strictly positive singular value,
ensuring numerical stability when evaluating sparse matrices where $W = \hat{W} \odot Z$.

Feng et al. \cite{lipschitz_condition_number} examined the effect of sparsity on the condition number and the Lipschitz constant of the network to capture
the resulting robustness of the sparse model. They note that highly pruned models tends to have a lower upper bound of the Lipschitz constant. However, sparse
weight matrices are more likely to be ill-conditioned, which is signaled by a high condition number.

By tracking these two metrics during the gradual pruning process, we can compare the effects of the aforementioned design decisions
on the robustness of the resulting sparse model, besides the standard accuracy metric and the length of the process.


